\documentclass[11pt]{article}
\usepackage[left=1in, right=1in, bottom=1in, top=1.5in]{geometry}

\begin{document}

\title{Why synecdols?}
\author{Juan A. Guarín-Rojas}
\date{\today}
\maketitle

The search of exact solutions improves the accuracy of our models of nature and the predictions derived from them. Among the exact solutions to the Einstein Field Equations, the solutions of an object with a quadrupole moment of mass (or higher) are desirable. This is because they admit a perturbative solution that emits gravitational waves. A solution like this involve an object with a geometry far away from the sphere like ellipsoids and cycloids. Another non-spheric surface is the Cartesian ovoid, which is defined as the locus of points such that its distances to two foci $r$ and $r'$ satisfy the relation $r+ar'=\rm{constant}$. The orthogonal surfaces to the ovoids allows to construct an orthogonal curvilinear system with ovoidal symmetry. Such surfaces are called synecdoloids inspired in the figure of speech synecdoche; and with the purpose to preserve the literature analogies between some conic curves and figures of speech like: parabola and parable, and hyperbola and hyperbole.

\textbf{Synecdoles}

\begin{itemize}
	\item The exact solutions of a non-spheric object are desirable because they have a quadrupole momentum of mass.
	\item The Cartesian ovoids are one generalization of the sphere and the ellipsoids that allows to construc a solution to the Einstein Field Equations (EE).
	\item The synecdols are the curves that together with the ovals form an orthogonal curvilinear coordinate system, in which the solutions to EE can be found.
\end{itemize}

\end{document}